\section{Discussion}

\noindent\textbf{Rethinking the safety calibration target.}
All evaluated mechanisms are calibrated on binary distinctions, concentrating protection at L3 while leaving L1--L2 underserved and steepening the severity gradient; deployments in sensitive contexts must evaluate L1--L2 explicitly rather than trust L3-derived claims.

\noindent\textbf{The cross-modal gap.}
Text filters cannot decode visually adversarial prompts, and image classifiers trained on explicit L3 content fail at L1--L2 violations (72.32\% of L3 prompts carry no explicit harmful keyword).

\noindent\textbf{From diagnosis to mitigation.}
Insight 6 turns the benchmark from a measurement tool into an instrument: fine-tuning Qwen2.5-VL-7B with LoRA on \dataname's own reference images and per-rung risk explanations (\SuppSec{\ref{app:remediation}}) recovers the L1--L2 band that zero-shot filters miss and transfers to fresh generations from all six evaluated models (FBR 2.4\% pooled, per-generator 2.2--2.7\%, under prompt-aware inference), and the recipe is taxonomy-agnostic (regenerate ladders under one's own policy, retrain).
Two attribution caveats remain: we do not ablate binary-labeled supervision on the same images, so whether the ladder \emph{structure} is causal or any boundary-dense labels suffice is untested (a binary-tuned control is planned for the extended analysis); and inference is description-conditioned, so deployment-realistic (description-free) operation of the tuned filter is unresolved.

\noindent\textbf{Scope boundaries.}
Copyright and PID exhibit structurally lower HGR than NSFW categories (\Tab{\ref{tab:t2i_models}}): MLLM classifiers are weaker at fine-grained identity harms (e.g., character likeness), so the safety gap in these categories is likely \emph{underestimated}.
Severity definitions also reflect a Western (US/EU) legal framework inherited from current models' safety-alignment data; a bounded non-Western probe of Public Figures and Copyright in \SuppSec{\ref{app:cultural}} quantifies the resulting asymmetry. We treat this as an explicit scope limitation and a foundation for a multilingual, multi-jurisdictional extension; a future release will add finer severity gradations at drastic L2$\to$L3 transitions.
